\section{Related Works}
\label{related_works}

Methods for extracting itemsets from data have been extensively studied in the field of itemset mining \cite{fournier2017survey}. Various pattern evaluation methods and extraction techniques have been proposed based on the specific characteristics of the itemsets. Recent trends in itemset and pattern mining are summarized in \cite{fournier2022pattern} and \cite{fournier2019survey}.

A representative task in this domain is frequent itemset mining, which involves extracting itemsets that appear often in the dataset. Methods such as Apriori \cite{agrawal1994fast}, ECLAT \cite{zaki1997new}, and FP‑Growth \cite{han2000mining} have been developed for this purpose. In contrast, dense patterns are itemsets that exhibit high frequency within specific intervals despite having low overall support. The aforementioned methods evaluate only global frequency and fail to capture local variations. As a result, a high minimum support threshold may overlook low‑frequency dense patterns, while a low threshold expands the search space with many non‑dense patterns—a dilemma known as the rare item problem \cite{liu1999mining}.

To overcome these challenges, research has focused on extracting patterns based on variations in frequency over time. For example, TP‑Mine \cite{wan2009discovering} extends an earlier method \cite{dong1999efficient} by splitting the data at a specified point and extracting itemsets with significant frequency changes between the resulting halves. However, the extracted patterns depend heavily on the chosen split point. SPF \cite{chang2002mining} calculates an itemset’s frequency over the span from its first to its last occurrence to identify high‑frequency segments; however, it cannot capture finer‑grained frequency changes within that span.
% limiting its ability to extract patterns that appear briefly but intensely.

As alternatives, methods based on the gaps between itemset occurrences have been proposed \cite{radhakrishna2015survey}. Periodic Frequent Pattern Mining(PFPM) \cite{tanbeer2009discovering} extracts itemsets that occur frequently overall and whose every gap between successive occurrences does not exceed a preset maximum gap. Although a short maximum gap can define density, it excludes dense patterns with variable spacing. Partial Periodic Frequent Pattern Mining(PPFPM) \cite{kiran2017discovering} relaxes this constraint by introducing a threshold on the ratio of gaps that meet the maximum gap. However, since it does not require these gaps to be consecutive, it cannot reliably distinguish truly dense intervals from sparse ones. Additionally, methods such as Locally Periodic Frequent Itemset Mining(LPFIM) \cite{mahanta2005finding} and Local Periodic Pattern Mining(LPPM) \cite{fournier2021mining} search for runs of occurrences in which each gap is below the maximum and the run length satisfies a minimum requirement. LPFIM, for example, identifies each contiguous run of occurrences whose gaps all lie beneath the threshold and whose length meets a minimum; its accuracy, however, is highly sensitive to the chosen gap threshold: if set too high, intervals become too long and dilute the local density; if too low, they fragment and fail to capture true bursts. Similarly, LPPM accumulates the difference between each gap and the maximum gap until a threshold is exceeded, but a high allowance admits non‑dense intervals and lowers precision. Overall, these gap‑based methods struggle to handle diverse periodicities while maintaining high extraction accuracy.

Swarm intelligence \cite{kennedy2006swarm} has been applied to reduce the computational cost of pattern evaluation through localized search. Swarm intelligence involves multiple agents cooperating to solve optimization problems and is embodied in algorithms such as Particle Swarm Optimization (PSO) \cite{kennedy1995particle} and Genetic Algorithms (GA) \cite{holland1975adaptation}. In these approaches, many agents explore the solution space, and the search narrows based on provisional optimal solutions, enabling rapid convergence to near‑optimal results. Pattern mining can be viewed as an optimization problem where the evaluation metric serves as the objective function and the extracted patterns are the solutions. Consequently, swarm intelligence techniques have been applied to pattern mining \cite{djenouri2018new}\cite{telikani2020survey}. Extensions of PSO for pattern mining, such as PSO‑FIM \cite{kuo2011application} and PSO‑Apriori \cite{djenouri2017combining}, represent each particle’s position as an itemset and update these positions based on evaluation values. Similarly, GA‑based methods such as AGA \cite{alatacs2006efficient}, GA‑FIM \cite{djenouri2014association}, and GA‑Apriori \cite{djenouri2017combining} evolve itemset representations through selection, crossover, and mutation. Although these swarm intelligence methods reduce the number of candidate patterns, they still require scanning the entire dataset for evaluation, which limits the overall efficiency.

The proposed approach distinguishes itself from existing works in two key aspects. First, it redefines dense patterns by removing any gap constraint from their definition of support, thereby eliminating difficult gap‑threshold tuning and enabling more accurate extraction. Second, it incorporates an agent‑based local search that focuses on dense intervals without scanning the entire dataset. By restricting evaluation to these spans, we can significantly reduces computational costs while preserving the accuracy.
